\subsection{Normal Primary Structure}

\begin{corollary} \label{cor:normal-kernels-of-powers-and-squarefree-minimal-polynomial} \hlt{Normal Operators Have No Nilpotent Primary Part}\\
Let $\mathcal{V}\neq\{\mathbf{0}\}$ be a finite-dimensional complex inner-product space and let $T\in\LL(\mathcal{V})$ be normal.
\begin{enumerate}[label=(\roman*)]
\setlength{\itemsep}{0pt}
\item $(\forall\,p\in\mathbb{C}[z])(\forall\,k\in\mathbb{Z}_{>0})\,\ker\!\bigl(p(T)^k\bigr)=\ker\!\bigl(p(T)\bigr)$.
\item $\mu_T(z)=\prod\limits_{\lambda\in\operatorname{spec}(T)}(z-\lambda)$, so $\mu_T$ is square-free.
\end{enumerate}
\begin{proof}
Set $n\defeq\dim(\mathcal{V})$ and $J\defeq\{1,\ldots,n\}$. By Theorem \ref{thm:complex-spectral-theorem}(i),(ii), choose an orthonormal eigenbasis $\beta=(\mathbf{e}_1,\ldots,\mathbf{e}_n)$.\\
Fix $(\lambda_1,\ldots,\lambda_n)\in\mathbb{C}^n$ satisfying $(\forall\,j\in J)\,T\mathbf{e}_j=\lambda_j\mathbf{e}_j$.
\begin{enumerate}[label=(\roman*)]
\setlength{\itemsep}{0pt}
\item Fix $p\in\mathbb{C}[z]$ and $k\in\mathbb{Z}_{>0}$, and set $S\defeq p(T)$. By Proposition \ref{prop:polynomial-eigenvector-evaluation}, $(\forall\,j\in J)\,S\mathbf{e}_j=p(\lambda_j)\mathbf{e}_j$.\\
The identity $(\forall\,j\in J)\,S^k\mathbf{e}_j=p(\lambda_j)^k\mathbf{e}_j$ holds by induction on $k$.\\
Hence $\ker(S^k)=\operatorname{span}\{\mathbf{e}_j:j\in J,\ p(\lambda_j)^k=0\}=\operatorname{span}\{\mathbf{e}_j:j\in J,\ p(\lambda_j)=0\}=\ker(S)$ by uniqueness of $\beta$-coordinates and $k>0$.
\item Set $A\defeq[T]_{\beta\leftarrow\beta}=\operatorname{diag}(\lambda_1,\ldots,\lambda_n)$.\\
Its distinct eigenvalues are precisely the elements of $\operatorname{spec}(T)$ by Proposition \ref{prop:matrix-representation-action} and Proposition \ref{prop:coordinate-map-isomorphism}(ii),(iii).\\
The identity $\mu_A(z)=\prod\limits_{\lambda\in\operatorname{spec}(T)}(z-\lambda)$ holds by Theorem \ref{thm:diagonalizable-iff-minimal-polynomial-squarefree}(i),(ii).\\
The identity $(\forall\,f\in\mathbb{C}[z])\,[f(T)]_{\beta\leftarrow\beta}=f(A)$ holds by Proposition \ref{prop:matrix-representation-respects-linear-combinations} and Proposition \ref{prop:composition-and-matrix-multiplication}.\\
Thus $(\forall\,f\in\mathbb{C}[z])\,[f(T)=0\Longleftrightarrow f(A)=0]$ by Theorem \ref{thm:linear-maps-and-matrices-same-data}. Hence $\mu_T=\mu_A$ by Definitions \ref{def:operator-minimal-polynomial} and \ref{def:minimal-polynomial-matrix}.\\
Therefore $\mu_T(z)=\prod\limits_{\lambda\in\operatorname{spec}(T)}(z-\lambda)$. Its linear factors are pairwise distinct, so $\mu_T$ is square-free.
\end{enumerate}
\hlt{Structural Commentary}: The spectral basis controls polynomial kernels by scalar evaluation and identifies $\mu_T$ by diagonalizability.\\
\hlt{Key Insight}: Positive powers preserve the zero set of each scalar $p(\lambda)$.
\end{proof}
\end{corollary}

\begin{proposition} \label{prop:normal-coprime-invariant-subspaces-orthogonal} \hlt{Coprime Normal Primary Pieces Are Orthogonal}\\
Let $\mathcal{V}$ be a finite-dimensional complex inner-product space, let $T\in\LL(\mathcal{V})$ be normal, and let $\mathcal{S},\mathcal{R}\leq\mathcal{V}$ be $T$-invariant.\\
Let $(p,q,a,b)\in\mathbb{C}[z]^4$ satisfy $ap+bq=1$, $p(T)|_{\mathcal{S}}=0$, and $q(T)|_{\mathcal{R}}=0$. Then $\mathcal{S}\perp\mathcal{R}$.
\begin{proof}
Fix $\mathbf{s}\in\mathcal{S}$ and $\mathbf{r}\in\mathcal{R}$.
\begin{enumerate}[label=(\roman*)]
\setlength{\itemsep}{0pt}
\item \hlt{Bézout Reduction}: By Proposition \ref{prop:operator-polynomial-products}(i) and $ap+bq=1$, $I_{\mathcal{V}}=a(T)p(T)+b(T)q(T)$. Hence $\mathbf{s}=b(T)q(T)\mathbf{s}$ by $p(T)|_{\mathcal{S}}=0$.
\item \hlt{Adjoint Annihilation}: By Corollary \ref{cor:normal-operator-invariant-subspaces-reduce}(ii) and Proposition \ref{prop:adjoints-inverses-and-polynomials}(ii), $b(T)^*\mathcal{R},q(T)^*\mathcal{R}\subseteq\mathcal{R}$.\\
For $\mathbf{u},\mathbf{v}\in\mathcal{R}$, $\innerprod{q(T)^*\mathbf{u},\mathbf{v}}=\innerprod{\mathbf{u},q(T)\mathbf{v}}=0$ by $q(T)|_{\mathcal{R}}=0$. Hence $q(T)^*\mathcal{R}\subseteq\mathcal{R}\cap\mathcal{R}^{\perp}=\{\mathbf{0}\}$, so $q(T)^*|_{\mathcal{R}}=0$.
\item \hlt{Orthogonality}: Since $b(T)^*\mathbf{r}\in\mathcal{R}$, $\innerprod{\mathbf{s},\mathbf{r}}=\innerprod{b(T)q(T)\mathbf{s},\mathbf{r}}=\innerprod{\mathbf{s},q(T)^*b(T)^*\mathbf{r}}=0$. Since $\mathbf{s}$ and $\mathbf{r}$ were arbitrary, $\mathcal{S}\perp\mathcal{R}$.
\end{enumerate}
\hlt{Structural Commentary}: The Bézout identity transfers $\mathbf{s}$ to $q(T)$, while normality transfers $q(T)$-annihilation to $q(T)^*$.\\
\hlt{Key Insight}: Adjoint invariance and $q(T)|_{\mathcal{R}}=0$ place $q(T)^*\mathcal{R}$ in $\mathcal{R}\cap\mathcal{R}^{\perp}$.
\end{proof}
\end{proposition}
